\section{Reverse a Doubly Linked List}
\label{sec:ll-reverse-dll}

\problemstatement{Reverse a doubly linked list in place so its head becomes
the tail and vice versa, returning the new head.}

\begin{patternbox}
Because each DLL node already knows both neighbours, reversing is just
\textbf{swapping every node's \texttt{prev} and \texttt{next}} pointers.
After swapping all of them, the old tail -- whose \texttt{next} is now null
-- is the new head.
\end{patternbox}

\subsection*{Swap the two pointers at every node}

Walk the list; at each node exchange its \texttt{prev} and \texttt{next}.
Once swapped, to advance you must follow the \emph{old} \texttt{next}, which
is now stored in \texttt{prev} -- so move via the just-swapped
\texttt{prev}. When you fall off the end, the last node you swapped is the
new head. No new nodes are created; only pointers are flipped.

\begin{ideabox}[Reversal Is Local Pointer Swapping]
In a DLL, the global reversal decomposes into an identical local operation
at each node: swap its two links. Because the links are symmetric, doing
this everywhere reverses the whole list without tracking a running
\texttt{prev} as a singly linked reversal must.
\end{ideabox}

\subsection*{C++ implementation}

\begin{lstlisting}
#include <bits/stdc++.h>
using namespace std;

struct Node {
    int val;
    Node* prev;
    Node* next;
    Node(int v) : val(v), prev(nullptr), next(nullptr) {}
};

Node* build(const vector<int>& a) {
    Node dummy(0), *tail = &dummy;
    for (int x : a) { Node* n = new Node(x); tail->next = n; n->prev = tail; tail = n; }
    Node* head = dummy.next; if (head) head->prev = nullptr;
    return head;
}

Node* reverse(Node* head) {
    Node* curr = head;
    Node* last = head;
    while (curr) {
        last = curr;
        swap(curr->prev, curr->next);             // flip both links
        curr = curr->prev;                        // old next is now in prev
    }
    return last;                                  // old tail = new head
}

void print(Node* head) {
    for (Node* c = head; c; c = c->next) cout << c->val << " ";
    cout << "\n";
}

int main() {
    Node* head = build({1, 2, 3, 4});
    head = reverse(head);
    print(head);                                  // 4 3 2 1
    return 0;
}
\end{lstlisting}

\begin{complexitybox}
\textbf{Time Complexity.} $O(n)$ -- one pass. \textbf{Space Complexity.}
$O(1)$ -- only pointer swaps.
\end{complexitybox}

\begin{takeawaybox}
Reversing a DLL is swapping \texttt{prev} and \texttt{next} at every node --
the symmetric structure makes global reversal a local operation. Contrast
with the singly linked reversal (next section) which must thread a running
\texttt{prev} because there is no backward pointer.
\end{takeawaybox}
